% =============================================================================
%  Fundamentos compartilhados --- Módulo 02: linguagens de programação
%  Incluído pelas trilhas Kotlin e Flutter com % =============================================================================
%  Fundamentos compartilhados --- Módulo 02: linguagens de programação
%  Incluído pelas trilhas Kotlin e Flutter com % =============================================================================
%  Fundamentos compartilhados --- Módulo 02: linguagens de programação
%  Incluído pelas trilhas Kotlin e Flutter com % =============================================================================
%  Fundamentos compartilhados --- Módulo 02: linguagens de programação
%  Incluído pelas trilhas Kotlin e Flutter com \input{../fundamentos/f02-linguagens}
% =============================================================================
\section{Fundamentos: linguagens de programação}

\begin{frame}{Paradigmas de programação}
  \begin{columns}[T]
    \begin{column}{0.54\textwidth}
      \scriptsize
      \begin{tabularx}{\linewidth}{@{}l L@{}}
        \toprule
        \cabecalho{Paradigma} & \cabecalho{Ideia central} \\
        \midrule
        Imperativo & sequência de comandos que alteram o estado (\emph{faça isto, depois aquilo}) \\
        Orientado a objetos & dados e comportamentos agrupados em objetos que trocam mensagens \\
        Funcional & funções puras, valores imutáveis, composição de funções \\
        Declarativo & descreve \textbf{o que} se quer, não \textbf{como} obter (SQL, interfaces em Compose e Flutter) \\
        \bottomrule
      \end{tabularx}
    \end{column}
    \begin{column}{0.44\textwidth}
      \begin{definicao}[Paradigma]
        Estilo de organizar o raciocínio e o código. Linguagens modernas são \textbf{multiparadigma}: Kotlin e Dart misturam imperativo, objetos e funcional.
      \end{definicao}
      \begin{dica}
        Use o paradigma que deixa cada parte mais clara: objetos para modelar o domínio, funções para transformar dados, declarativo para descrever a interface.
      \end{dica}
    \end{column}
  \end{columns}
  \begin{historia}
    Fortran (1957) era imperativo; Lisp (1958), funcional; Simula (1967) e Smalltalk (anos 1970) criaram a orientação a objetos. Kotlin (2011, JetBrains) e Dart (2011, Google) herdaram ideias de todos.
  \end{historia}
\end{frame}

\begin{frame}{Sistemas de tipos}
  \begin{columns}[T]
    \begin{column}{0.5\textwidth}
      \begin{definicao}[Tipo]
        Conjunto de valores possíveis e das operações permitidas sobre eles. Somar dois \cod{Int} faz sentido; somar um \cod{Int} a uma \cod{Bancada}, não.
      \end{definicao}
      \scriptsize
      \begin{tabularx}{\linewidth}{@{}l L@{}}
        \toprule
        \cabecalho{Classificação} & \cabecalho{Significado} \\
        \midrule
        Estática & tipos verificados na compilação (Kotlin, Dart, Java) \\
        Dinâmica & tipos verificados na execução (Python, JavaScript) \\
        Forte & conversões perigosas são proibidas ou explícitas \\
        Inferência & o compilador deduz o tipo (\cod{val x = 3}) \\
        \bottomrule
      \end{tabularx}
    \end{column}
    \begin{column}{0.48\textwidth}
      \begin{conceito}[Por que tipagem estática ajuda]
        Muitos erros aparecem \textbf{antes} de o app rodar; a IDE completa código e refatora com segurança; o tipo documenta a intenção.
      \end{conceito}
      \begin{pordentro}[Solidez (\emph{soundness})]
        Um sistema de tipos é sólido quando garante que, se o programa compila, uma variável nunca terá em execução um valor de outro tipo. O Dart 3 é sólido, inclusive quanto a nulos.
      \end{pordentro}
    \end{column}
  \end{columns}
\end{frame}

\begin{frame}{O ``erro de um bilhão de dólares'': o nulo}
  \begin{columns}[T]
    \begin{column}{0.5\textwidth}
      \begin{historia}
        Tony Hoare introduziu a referência nula na linguagem ALGOL W em 1965. Décadas depois, chamou-a de seu ``erro de um bilhão de dólares'', pelo custo das falhas que ela causou.
      \end{historia}
      \begin{definicao}[\emph{Null safety}]
        O sistema de tipos separa tipos que \textbf{podem} ser nulos (\cod{String?}) dos que \textbf{não podem} (\cod{String}). Usar um valor possivelmente nulo sem tratar vira erro de compilação, não falha em produção.
      \end{definicao}
    \end{column}
    \begin{column}{0.48\textwidth}
      \centering
      \begin{tikzpicture}[node distance=3mm,every node/.style={font=\scriptsize,align=center}]
        \node[caixa,minimum width=2.2cm] (a) {\cod{String?}\\ \tiny pode ser nulo};
        \node[caixapetroleo,right=8mm of a,minimum width=2.2cm] (b) {\cod{String}\\ \tiny nunca nulo};
        \draw[seta] (a) -- node[rotulo,above]{\tiny verificação} (b);
        \node[below=3mm of a,font=\tiny,text width=2.2cm,align=center] {\cod{if (x != null)}\\ \cod{x ?: padrão}};
      \end{tikzpicture}
      \begin{pordentro}[Análise de fluxo]
        Depois de um \cod{if (x != null)}, o compilador sabe que \cod{x} não é nulo naquele trecho e permite usá-lo diretamente (\emph{smart cast} no Kotlin, promoção de tipo no Dart).
      \end{pordentro}
      \begin{atencao}
        O operador \cod{!!} (Kotlin) ou \cod{!} (Dart) desliga a proteção: só use quando tiver certeza.
      \end{atencao}
    \end{column}
  \end{columns}
\end{frame}

\begin{frame}{Os pilares da orientação a objetos}
  \begin{columns}[T]
    \begin{column}{0.54\textwidth}
      \scriptsize
      \begin{tabularx}{\linewidth}{@{}l L@{}}
        \toprule
        \cabecalho{Pilar} & \cabecalho{Significado e exemplo} \\
        \midrule
        Encapsulamento & esconder o estado interno e expor operações: a \cod{Reserva} só muda pelo método \cod{cancelar()} \\
        Abstração & mostrar o essencial e ocultar detalhes: \cod{FonteLeituras} sem dizer se é BLE ou MQTT \\
        Herança & uma classe especializa outra: \cod{SensorTemperatura} herda de \cod{Sensor} \\
        Polimorfismo & o mesmo comando, comportamentos diferentes: \cod{fonte.leituras()} em qualquer implementação \\
        \bottomrule
      \end{tabularx}
    \end{column}
    \begin{column}{0.44\textwidth}
      \begin{definicao}[Objeto e classe]
        A \textbf{classe} é o molde (atributos e métodos); o \textbf{objeto} é uma instância concreta, com estado próprio.
      \end{definicao}
      \begin{dica}[Composição em vez de herança]
        Prefira montar objetos com outros objetos (``tem um'') a criar hierarquias profundas (``é um''): mudanças ficam localizadas e o código, mais flexível.
      \end{dica}
    \end{column}
  \end{columns}
  \begin{analogia}
    Uma tomada é uma abstração: qualquer aparelho com o plugue certo funciona (polimorfismo), e ninguém precisa conhecer a fiação dentro da parede (encapsulamento).
  \end{analogia}
\end{frame}

\begin{frame}{Programação funcional}
  \begin{columns}[T]
    \begin{column}{0.5\textwidth}
      \begin{definicao}[Função pura]
        Para a mesma entrada, devolve sempre a mesma saída e não causa efeitos colaterais (não altera variáveis externas, não grava em disco, não acessa a rede).
      \end{definicao}
      \begin{definicao}[Funções de primeira classe e de ordem superior]
        Funções são valores: podem ser guardadas em variáveis, passadas como argumento e devolvidas. Uma função que recebe ou devolve outra é de \textbf{ordem superior} (\cod{map}, \cod{filter}).
      \end{definicao}
    \end{column}
    \begin{column}{0.48\textwidth}
      \centering
      \begin{tikzpicture}[node distance=4mm,every node/.style={font=\scriptsize,align=center}]
        \node[caixa,minimum width=1.9cm] (l) {leituras};
        \node[caixapetroleo,right=of l,minimum width=1.6cm] (f) {\cod{filter}\\ \tiny > 30\,°C};
        \node[caixapetroleo,right=of f,minimum width=1.6cm] (m) {\cod{map}\\ \tiny para alerta};
        \node[caixaambar,below=4mm of m,minimum width=1.6cm] (r) {alertas};
        \draw[seta] (l) -- (f); \draw[seta] (f) -- (m); \draw[seta] (m) -- (r);
      \end{tikzpicture}
      \begin{conceito}[Por que importa no mobile]
        Funções puras são fáceis de testar e seguras entre \emph{threads}; transformar listas com \cod{map}/\cod{filter} é mais legível que laços; interfaces declarativas tratam a tela como uma função do estado (Módulo 05).
      \end{conceito}
    \end{column}
  \end{columns}
\end{frame}

\begin{frame}{Memória: pilha, \emph{heap} e coleta de lixo}
  \begin{columns}[c]
    \begin{column}{0.46\textwidth}
      \centering
      \begin{tikzpicture}[every node/.style={font=\tiny}]
        \node[font=\scriptsize\bfseries,text=azulpetroleo] at (0.9,2.7) {Pilha};
        \foreach \y/\t in {2.2/main(),1.7/carregar(),1.2/calcular()} \node[caixa,minimum width=1.8cm] at (0.9,\y) {\t};
        \node[font=\scriptsize\bfseries,text=azulpetroleo] at (3.7,2.7) {\emph{Heap}};
        \draw[cinzamedio,rounded corners] (2.6,0.3) rectangle (4.8,2.45);
        \node[caixapetroleo,minimum width=1.3cm] (o1) at (3.3,1.95) {Bancada};
        \node[caixapetroleo,minimum width=1.3cm] (o2) at (4.1,1.2) {Lista};
        \node[caixa,minimum width=1.3cm,fill=cinzaclaro] (o3) at (3.3,0.65) {sem ref.};
        \draw[seta] (1.8,2.2) -- (o1); \draw[seta] (1.8,1.7) -- (o2); \draw[seta] (o2) -- (o1);
        \node[font=\tiny,text=vermelhoerro] at (3.7,0.2) {coletado};
      \end{tikzpicture}
    \end{column}
    \begin{column}{0.52\textwidth}
      \begin{definicao}[Pilha e \emph{heap}]
        A \textbf{pilha} guarda as chamadas de função e suas variáveis locais, liberadas ao retornar. O \textbf{\emph{heap}} guarda os objetos, que vivem enquanto houver referências a eles.
      \end{definicao}
      \begin{definicao}[Coleta de lixo]
        O \emph{runtime} encontra objetos que não são mais alcançáveis a partir das referências vivas e libera sua memória automaticamente.
      \end{definicao}
      \begin{pordentro}[Gerações]
        A maioria dos objetos morre jovem. Por isso o ART e a VM do Dart separam objetos novos e antigos e coletam os novos com frequência e rapidez.
      \end{pordentro}
    \end{column}
  \end{columns}
\end{frame}

\begin{frame}{Imutabilidade e igualdade}
  \begin{columns}[T]
    \begin{column}{0.5\textwidth}
      \begin{definicao}[Imutável]
        Objeto que não muda depois de criado. Para ``alterar'', cria-se uma cópia com o campo novo (\cod{copy} no Kotlin, \cod{copyWith} no Dart).
      \end{definicao}
      \begin{conceito}[Vantagens]
        \begin{itemize}
          \item Ninguém altera o objeto ``pelas costas''.
          \item Seguro para compartilhar entre \emph{threads} e \emph{isolates}.
          \item Comparar o antes e o depois de um estado fica trivial (Módulo 07).
        \end{itemize}
      \end{conceito}
    \end{column}
    \begin{column}{0.48\textwidth}
      \begin{definicao}[Igualdade de valor e de identidade]
        \textbf{Identidade}: as duas variáveis apontam para o \textbf{mesmo} objeto. \textbf{Valor}: objetos diferentes com o mesmo conteúdo. \cod{data class} (Kotlin) e \emph{records} (Dart) comparam por valor.
      \end{definicao}
      \begin{erro}[Lista mutável compartilhada]
        Dois componentes recebem a mesma lista; um a altera e o outro mostra dados que não esperava. Exponha listas somente leitura.
      \end{erro}
    \end{column}
  \end{columns}
\end{frame}

% =============================================================================
\section{Fundamentos: linguagens de programação}

\begin{frame}{Paradigmas de programação}
  \begin{columns}[T]
    \begin{column}{0.54\textwidth}
      \scriptsize
      \begin{tabularx}{\linewidth}{@{}l L@{}}
        \toprule
        \cabecalho{Paradigma} & \cabecalho{Ideia central} \\
        \midrule
        Imperativo & sequência de comandos que alteram o estado (\emph{faça isto, depois aquilo}) \\
        Orientado a objetos & dados e comportamentos agrupados em objetos que trocam mensagens \\
        Funcional & funções puras, valores imutáveis, composição de funções \\
        Declarativo & descreve \textbf{o que} se quer, não \textbf{como} obter (SQL, interfaces em Compose e Flutter) \\
        \bottomrule
      \end{tabularx}
    \end{column}
    \begin{column}{0.44\textwidth}
      \begin{definicao}[Paradigma]
        Estilo de organizar o raciocínio e o código. Linguagens modernas são \textbf{multiparadigma}: Kotlin e Dart misturam imperativo, objetos e funcional.
      \end{definicao}
      \begin{dica}
        Use o paradigma que deixa cada parte mais clara: objetos para modelar o domínio, funções para transformar dados, declarativo para descrever a interface.
      \end{dica}
    \end{column}
  \end{columns}
  \begin{historia}
    Fortran (1957) era imperativo; Lisp (1958), funcional; Simula (1967) e Smalltalk (anos 1970) criaram a orientação a objetos. Kotlin (2011, JetBrains) e Dart (2011, Google) herdaram ideias de todos.
  \end{historia}
\end{frame}

\begin{frame}{Sistemas de tipos}
  \begin{columns}[T]
    \begin{column}{0.5\textwidth}
      \begin{definicao}[Tipo]
        Conjunto de valores possíveis e das operações permitidas sobre eles. Somar dois \cod{Int} faz sentido; somar um \cod{Int} a uma \cod{Bancada}, não.
      \end{definicao}
      \scriptsize
      \begin{tabularx}{\linewidth}{@{}l L@{}}
        \toprule
        \cabecalho{Classificação} & \cabecalho{Significado} \\
        \midrule
        Estática & tipos verificados na compilação (Kotlin, Dart, Java) \\
        Dinâmica & tipos verificados na execução (Python, JavaScript) \\
        Forte & conversões perigosas são proibidas ou explícitas \\
        Inferência & o compilador deduz o tipo (\cod{val x = 3}) \\
        \bottomrule
      \end{tabularx}
    \end{column}
    \begin{column}{0.48\textwidth}
      \begin{conceito}[Por que tipagem estática ajuda]
        Muitos erros aparecem \textbf{antes} de o app rodar; a IDE completa código e refatora com segurança; o tipo documenta a intenção.
      \end{conceito}
      \begin{pordentro}[Solidez (\emph{soundness})]
        Um sistema de tipos é sólido quando garante que, se o programa compila, uma variável nunca terá em execução um valor de outro tipo. O Dart 3 é sólido, inclusive quanto a nulos.
      \end{pordentro}
    \end{column}
  \end{columns}
\end{frame}

\begin{frame}{O ``erro de um bilhão de dólares'': o nulo}
  \begin{columns}[T]
    \begin{column}{0.5\textwidth}
      \begin{historia}
        Tony Hoare introduziu a referência nula na linguagem ALGOL W em 1965. Décadas depois, chamou-a de seu ``erro de um bilhão de dólares'', pelo custo das falhas que ela causou.
      \end{historia}
      \begin{definicao}[\emph{Null safety}]
        O sistema de tipos separa tipos que \textbf{podem} ser nulos (\cod{String?}) dos que \textbf{não podem} (\cod{String}). Usar um valor possivelmente nulo sem tratar vira erro de compilação, não falha em produção.
      \end{definicao}
    \end{column}
    \begin{column}{0.48\textwidth}
      \centering
      \begin{tikzpicture}[node distance=3mm,every node/.style={font=\scriptsize,align=center}]
        \node[caixa,minimum width=2.2cm] (a) {\cod{String?}\\ \tiny pode ser nulo};
        \node[caixapetroleo,right=8mm of a,minimum width=2.2cm] (b) {\cod{String}\\ \tiny nunca nulo};
        \draw[seta] (a) -- node[rotulo,above]{\tiny verificação} (b);
        \node[below=3mm of a,font=\tiny,text width=2.2cm,align=center] {\cod{if (x != null)}\\ \cod{x ?: padrão}};
      \end{tikzpicture}
      \begin{pordentro}[Análise de fluxo]
        Depois de um \cod{if (x != null)}, o compilador sabe que \cod{x} não é nulo naquele trecho e permite usá-lo diretamente (\emph{smart cast} no Kotlin, promoção de tipo no Dart).
      \end{pordentro}
      \begin{atencao}
        O operador \cod{!!} (Kotlin) ou \cod{!} (Dart) desliga a proteção: só use quando tiver certeza.
      \end{atencao}
    \end{column}
  \end{columns}
\end{frame}

\begin{frame}{Os pilares da orientação a objetos}
  \begin{columns}[T]
    \begin{column}{0.54\textwidth}
      \scriptsize
      \begin{tabularx}{\linewidth}{@{}l L@{}}
        \toprule
        \cabecalho{Pilar} & \cabecalho{Significado e exemplo} \\
        \midrule
        Encapsulamento & esconder o estado interno e expor operações: a \cod{Reserva} só muda pelo método \cod{cancelar()} \\
        Abstração & mostrar o essencial e ocultar detalhes: \cod{FonteLeituras} sem dizer se é BLE ou MQTT \\
        Herança & uma classe especializa outra: \cod{SensorTemperatura} herda de \cod{Sensor} \\
        Polimorfismo & o mesmo comando, comportamentos diferentes: \cod{fonte.leituras()} em qualquer implementação \\
        \bottomrule
      \end{tabularx}
    \end{column}
    \begin{column}{0.44\textwidth}
      \begin{definicao}[Objeto e classe]
        A \textbf{classe} é o molde (atributos e métodos); o \textbf{objeto} é uma instância concreta, com estado próprio.
      \end{definicao}
      \begin{dica}[Composição em vez de herança]
        Prefira montar objetos com outros objetos (``tem um'') a criar hierarquias profundas (``é um''): mudanças ficam localizadas e o código, mais flexível.
      \end{dica}
    \end{column}
  \end{columns}
  \begin{analogia}
    Uma tomada é uma abstração: qualquer aparelho com o plugue certo funciona (polimorfismo), e ninguém precisa conhecer a fiação dentro da parede (encapsulamento).
  \end{analogia}
\end{frame}

\begin{frame}{Programação funcional}
  \begin{columns}[T]
    \begin{column}{0.5\textwidth}
      \begin{definicao}[Função pura]
        Para a mesma entrada, devolve sempre a mesma saída e não causa efeitos colaterais (não altera variáveis externas, não grava em disco, não acessa a rede).
      \end{definicao}
      \begin{definicao}[Funções de primeira classe e de ordem superior]
        Funções são valores: podem ser guardadas em variáveis, passadas como argumento e devolvidas. Uma função que recebe ou devolve outra é de \textbf{ordem superior} (\cod{map}, \cod{filter}).
      \end{definicao}
    \end{column}
    \begin{column}{0.48\textwidth}
      \centering
      \begin{tikzpicture}[node distance=4mm,every node/.style={font=\scriptsize,align=center}]
        \node[caixa,minimum width=1.9cm] (l) {leituras};
        \node[caixapetroleo,right=of l,minimum width=1.6cm] (f) {\cod{filter}\\ \tiny > 30\,°C};
        \node[caixapetroleo,right=of f,minimum width=1.6cm] (m) {\cod{map}\\ \tiny para alerta};
        \node[caixaambar,below=4mm of m,minimum width=1.6cm] (r) {alertas};
        \draw[seta] (l) -- (f); \draw[seta] (f) -- (m); \draw[seta] (m) -- (r);
      \end{tikzpicture}
      \begin{conceito}[Por que importa no mobile]
        Funções puras são fáceis de testar e seguras entre \emph{threads}; transformar listas com \cod{map}/\cod{filter} é mais legível que laços; interfaces declarativas tratam a tela como uma função do estado (Módulo 05).
      \end{conceito}
    \end{column}
  \end{columns}
\end{frame}

\begin{frame}{Memória: pilha, \emph{heap} e coleta de lixo}
  \begin{columns}[c]
    \begin{column}{0.46\textwidth}
      \centering
      \begin{tikzpicture}[every node/.style={font=\tiny}]
        \node[font=\scriptsize\bfseries,text=azulpetroleo] at (0.9,2.7) {Pilha};
        \foreach \y/\t in {2.2/main(),1.7/carregar(),1.2/calcular()} \node[caixa,minimum width=1.8cm] at (0.9,\y) {\t};
        \node[font=\scriptsize\bfseries,text=azulpetroleo] at (3.7,2.7) {\emph{Heap}};
        \draw[cinzamedio,rounded corners] (2.6,0.3) rectangle (4.8,2.45);
        \node[caixapetroleo,minimum width=1.3cm] (o1) at (3.3,1.95) {Bancada};
        \node[caixapetroleo,minimum width=1.3cm] (o2) at (4.1,1.2) {Lista};
        \node[caixa,minimum width=1.3cm,fill=cinzaclaro] (o3) at (3.3,0.65) {sem ref.};
        \draw[seta] (1.8,2.2) -- (o1); \draw[seta] (1.8,1.7) -- (o2); \draw[seta] (o2) -- (o1);
        \node[font=\tiny,text=vermelhoerro] at (3.7,0.2) {coletado};
      \end{tikzpicture}
    \end{column}
    \begin{column}{0.52\textwidth}
      \begin{definicao}[Pilha e \emph{heap}]
        A \textbf{pilha} guarda as chamadas de função e suas variáveis locais, liberadas ao retornar. O \textbf{\emph{heap}} guarda os objetos, que vivem enquanto houver referências a eles.
      \end{definicao}
      \begin{definicao}[Coleta de lixo]
        O \emph{runtime} encontra objetos que não são mais alcançáveis a partir das referências vivas e libera sua memória automaticamente.
      \end{definicao}
      \begin{pordentro}[Gerações]
        A maioria dos objetos morre jovem. Por isso o ART e a VM do Dart separam objetos novos e antigos e coletam os novos com frequência e rapidez.
      \end{pordentro}
    \end{column}
  \end{columns}
\end{frame}

\begin{frame}{Imutabilidade e igualdade}
  \begin{columns}[T]
    \begin{column}{0.5\textwidth}
      \begin{definicao}[Imutável]
        Objeto que não muda depois de criado. Para ``alterar'', cria-se uma cópia com o campo novo (\cod{copy} no Kotlin, \cod{copyWith} no Dart).
      \end{definicao}
      \begin{conceito}[Vantagens]
        \begin{itemize}
          \item Ninguém altera o objeto ``pelas costas''.
          \item Seguro para compartilhar entre \emph{threads} e \emph{isolates}.
          \item Comparar o antes e o depois de um estado fica trivial (Módulo 07).
        \end{itemize}
      \end{conceito}
    \end{column}
    \begin{column}{0.48\textwidth}
      \begin{definicao}[Igualdade de valor e de identidade]
        \textbf{Identidade}: as duas variáveis apontam para o \textbf{mesmo} objeto. \textbf{Valor}: objetos diferentes com o mesmo conteúdo. \cod{data class} (Kotlin) e \emph{records} (Dart) comparam por valor.
      \end{definicao}
      \begin{erro}[Lista mutável compartilhada]
        Dois componentes recebem a mesma lista; um a altera e o outro mostra dados que não esperava. Exponha listas somente leitura.
      \end{erro}
    \end{column}
  \end{columns}
\end{frame}

% =============================================================================
\section{Fundamentos: linguagens de programação}

\begin{frame}{Paradigmas de programação}
  \begin{columns}[T]
    \begin{column}{0.54\textwidth}
      \scriptsize
      \begin{tabularx}{\linewidth}{@{}l L@{}}
        \toprule
        \cabecalho{Paradigma} & \cabecalho{Ideia central} \\
        \midrule
        Imperativo & sequência de comandos que alteram o estado (\emph{faça isto, depois aquilo}) \\
        Orientado a objetos & dados e comportamentos agrupados em objetos que trocam mensagens \\
        Funcional & funções puras, valores imutáveis, composição de funções \\
        Declarativo & descreve \textbf{o que} se quer, não \textbf{como} obter (SQL, interfaces em Compose e Flutter) \\
        \bottomrule
      \end{tabularx}
    \end{column}
    \begin{column}{0.44\textwidth}
      \begin{definicao}[Paradigma]
        Estilo de organizar o raciocínio e o código. Linguagens modernas são \textbf{multiparadigma}: Kotlin e Dart misturam imperativo, objetos e funcional.
      \end{definicao}
      \begin{dica}
        Use o paradigma que deixa cada parte mais clara: objetos para modelar o domínio, funções para transformar dados, declarativo para descrever a interface.
      \end{dica}
    \end{column}
  \end{columns}
  \begin{historia}
    Fortran (1957) era imperativo; Lisp (1958), funcional; Simula (1967) e Smalltalk (anos 1970) criaram a orientação a objetos. Kotlin (2011, JetBrains) e Dart (2011, Google) herdaram ideias de todos.
  \end{historia}
\end{frame}

\begin{frame}{Sistemas de tipos}
  \begin{columns}[T]
    \begin{column}{0.5\textwidth}
      \begin{definicao}[Tipo]
        Conjunto de valores possíveis e das operações permitidas sobre eles. Somar dois \cod{Int} faz sentido; somar um \cod{Int} a uma \cod{Bancada}, não.
      \end{definicao}
      \scriptsize
      \begin{tabularx}{\linewidth}{@{}l L@{}}
        \toprule
        \cabecalho{Classificação} & \cabecalho{Significado} \\
        \midrule
        Estática & tipos verificados na compilação (Kotlin, Dart, Java) \\
        Dinâmica & tipos verificados na execução (Python, JavaScript) \\
        Forte & conversões perigosas são proibidas ou explícitas \\
        Inferência & o compilador deduz o tipo (\cod{val x = 3}) \\
        \bottomrule
      \end{tabularx}
    \end{column}
    \begin{column}{0.48\textwidth}
      \begin{conceito}[Por que tipagem estática ajuda]
        Muitos erros aparecem \textbf{antes} de o app rodar; a IDE completa código e refatora com segurança; o tipo documenta a intenção.
      \end{conceito}
      \begin{pordentro}[Solidez (\emph{soundness})]
        Um sistema de tipos é sólido quando garante que, se o programa compila, uma variável nunca terá em execução um valor de outro tipo. O Dart 3 é sólido, inclusive quanto a nulos.
      \end{pordentro}
    \end{column}
  \end{columns}
\end{frame}

\begin{frame}{O ``erro de um bilhão de dólares'': o nulo}
  \begin{columns}[T]
    \begin{column}{0.5\textwidth}
      \begin{historia}
        Tony Hoare introduziu a referência nula na linguagem ALGOL W em 1965. Décadas depois, chamou-a de seu ``erro de um bilhão de dólares'', pelo custo das falhas que ela causou.
      \end{historia}
      \begin{definicao}[\emph{Null safety}]
        O sistema de tipos separa tipos que \textbf{podem} ser nulos (\cod{String?}) dos que \textbf{não podem} (\cod{String}). Usar um valor possivelmente nulo sem tratar vira erro de compilação, não falha em produção.
      \end{definicao}
    \end{column}
    \begin{column}{0.48\textwidth}
      \centering
      \begin{tikzpicture}[node distance=3mm,every node/.style={font=\scriptsize,align=center}]
        \node[caixa,minimum width=2.2cm] (a) {\cod{String?}\\ \tiny pode ser nulo};
        \node[caixapetroleo,right=8mm of a,minimum width=2.2cm] (b) {\cod{String}\\ \tiny nunca nulo};
        \draw[seta] (a) -- node[rotulo,above]{\tiny verificação} (b);
        \node[below=3mm of a,font=\tiny,text width=2.2cm,align=center] {\cod{if (x != null)}\\ \cod{x ?: padrão}};
      \end{tikzpicture}
      \begin{pordentro}[Análise de fluxo]
        Depois de um \cod{if (x != null)}, o compilador sabe que \cod{x} não é nulo naquele trecho e permite usá-lo diretamente (\emph{smart cast} no Kotlin, promoção de tipo no Dart).
      \end{pordentro}
      \begin{atencao}
        O operador \cod{!!} (Kotlin) ou \cod{!} (Dart) desliga a proteção: só use quando tiver certeza.
      \end{atencao}
    \end{column}
  \end{columns}
\end{frame}

\begin{frame}{Os pilares da orientação a objetos}
  \begin{columns}[T]
    \begin{column}{0.54\textwidth}
      \scriptsize
      \begin{tabularx}{\linewidth}{@{}l L@{}}
        \toprule
        \cabecalho{Pilar} & \cabecalho{Significado e exemplo} \\
        \midrule
        Encapsulamento & esconder o estado interno e expor operações: a \cod{Reserva} só muda pelo método \cod{cancelar()} \\
        Abstração & mostrar o essencial e ocultar detalhes: \cod{FonteLeituras} sem dizer se é BLE ou MQTT \\
        Herança & uma classe especializa outra: \cod{SensorTemperatura} herda de \cod{Sensor} \\
        Polimorfismo & o mesmo comando, comportamentos diferentes: \cod{fonte.leituras()} em qualquer implementação \\
        \bottomrule
      \end{tabularx}
    \end{column}
    \begin{column}{0.44\textwidth}
      \begin{definicao}[Objeto e classe]
        A \textbf{classe} é o molde (atributos e métodos); o \textbf{objeto} é uma instância concreta, com estado próprio.
      \end{definicao}
      \begin{dica}[Composição em vez de herança]
        Prefira montar objetos com outros objetos (``tem um'') a criar hierarquias profundas (``é um''): mudanças ficam localizadas e o código, mais flexível.
      \end{dica}
    \end{column}
  \end{columns}
  \begin{analogia}
    Uma tomada é uma abstração: qualquer aparelho com o plugue certo funciona (polimorfismo), e ninguém precisa conhecer a fiação dentro da parede (encapsulamento).
  \end{analogia}
\end{frame}

\begin{frame}{Programação funcional}
  \begin{columns}[T]
    \begin{column}{0.5\textwidth}
      \begin{definicao}[Função pura]
        Para a mesma entrada, devolve sempre a mesma saída e não causa efeitos colaterais (não altera variáveis externas, não grava em disco, não acessa a rede).
      \end{definicao}
      \begin{definicao}[Funções de primeira classe e de ordem superior]
        Funções são valores: podem ser guardadas em variáveis, passadas como argumento e devolvidas. Uma função que recebe ou devolve outra é de \textbf{ordem superior} (\cod{map}, \cod{filter}).
      \end{definicao}
    \end{column}
    \begin{column}{0.48\textwidth}
      \centering
      \begin{tikzpicture}[node distance=4mm,every node/.style={font=\scriptsize,align=center}]
        \node[caixa,minimum width=1.9cm] (l) {leituras};
        \node[caixapetroleo,right=of l,minimum width=1.6cm] (f) {\cod{filter}\\ \tiny > 30\,°C};
        \node[caixapetroleo,right=of f,minimum width=1.6cm] (m) {\cod{map}\\ \tiny para alerta};
        \node[caixaambar,below=4mm of m,minimum width=1.6cm] (r) {alertas};
        \draw[seta] (l) -- (f); \draw[seta] (f) -- (m); \draw[seta] (m) -- (r);
      \end{tikzpicture}
      \begin{conceito}[Por que importa no mobile]
        Funções puras são fáceis de testar e seguras entre \emph{threads}; transformar listas com \cod{map}/\cod{filter} é mais legível que laços; interfaces declarativas tratam a tela como uma função do estado (Módulo 05).
      \end{conceito}
    \end{column}
  \end{columns}
\end{frame}

\begin{frame}{Memória: pilha, \emph{heap} e coleta de lixo}
  \begin{columns}[c]
    \begin{column}{0.46\textwidth}
      \centering
      \begin{tikzpicture}[every node/.style={font=\tiny}]
        \node[font=\scriptsize\bfseries,text=azulpetroleo] at (0.9,2.7) {Pilha};
        \foreach \y/\t in {2.2/main(),1.7/carregar(),1.2/calcular()} \node[caixa,minimum width=1.8cm] at (0.9,\y) {\t};
        \node[font=\scriptsize\bfseries,text=azulpetroleo] at (3.7,2.7) {\emph{Heap}};
        \draw[cinzamedio,rounded corners] (2.6,0.3) rectangle (4.8,2.45);
        \node[caixapetroleo,minimum width=1.3cm] (o1) at (3.3,1.95) {Bancada};
        \node[caixapetroleo,minimum width=1.3cm] (o2) at (4.1,1.2) {Lista};
        \node[caixa,minimum width=1.3cm,fill=cinzaclaro] (o3) at (3.3,0.65) {sem ref.};
        \draw[seta] (1.8,2.2) -- (o1); \draw[seta] (1.8,1.7) -- (o2); \draw[seta] (o2) -- (o1);
        \node[font=\tiny,text=vermelhoerro] at (3.7,0.2) {coletado};
      \end{tikzpicture}
    \end{column}
    \begin{column}{0.52\textwidth}
      \begin{definicao}[Pilha e \emph{heap}]
        A \textbf{pilha} guarda as chamadas de função e suas variáveis locais, liberadas ao retornar. O \textbf{\emph{heap}} guarda os objetos, que vivem enquanto houver referências a eles.
      \end{definicao}
      \begin{definicao}[Coleta de lixo]
        O \emph{runtime} encontra objetos que não são mais alcançáveis a partir das referências vivas e libera sua memória automaticamente.
      \end{definicao}
      \begin{pordentro}[Gerações]
        A maioria dos objetos morre jovem. Por isso o ART e a VM do Dart separam objetos novos e antigos e coletam os novos com frequência e rapidez.
      \end{pordentro}
    \end{column}
  \end{columns}
\end{frame}

\begin{frame}{Imutabilidade e igualdade}
  \begin{columns}[T]
    \begin{column}{0.5\textwidth}
      \begin{definicao}[Imutável]
        Objeto que não muda depois de criado. Para ``alterar'', cria-se uma cópia com o campo novo (\cod{copy} no Kotlin, \cod{copyWith} no Dart).
      \end{definicao}
      \begin{conceito}[Vantagens]
        \begin{itemize}
          \item Ninguém altera o objeto ``pelas costas''.
          \item Seguro para compartilhar entre \emph{threads} e \emph{isolates}.
          \item Comparar o antes e o depois de um estado fica trivial (Módulo 07).
        \end{itemize}
      \end{conceito}
    \end{column}
    \begin{column}{0.48\textwidth}
      \begin{definicao}[Igualdade de valor e de identidade]
        \textbf{Identidade}: as duas variáveis apontam para o \textbf{mesmo} objeto. \textbf{Valor}: objetos diferentes com o mesmo conteúdo. \cod{data class} (Kotlin) e \emph{records} (Dart) comparam por valor.
      \end{definicao}
      \begin{erro}[Lista mutável compartilhada]
        Dois componentes recebem a mesma lista; um a altera e o outro mostra dados que não esperava. Exponha listas somente leitura.
      \end{erro}
    \end{column}
  \end{columns}
\end{frame}

% =============================================================================
\section{Fundamentos: linguagens de programação}

\begin{frame}{Paradigmas de programação}
  \begin{columns}[T]
    \begin{column}{0.54\textwidth}
      \scriptsize
      \begin{tabularx}{\linewidth}{@{}l L@{}}
        \toprule
        \cabecalho{Paradigma} & \cabecalho{Ideia central} \\
        \midrule
        Imperativo & sequência de comandos que alteram o estado (\emph{faça isto, depois aquilo}) \\
        Orientado a objetos & dados e comportamentos agrupados em objetos que trocam mensagens \\
        Funcional & funções puras, valores imutáveis, composição de funções \\
        Declarativo & descreve \textbf{o que} se quer, não \textbf{como} obter (SQL, interfaces em Compose e Flutter) \\
        \bottomrule
      \end{tabularx}
    \end{column}
    \begin{column}{0.44\textwidth}
      \begin{definicao}[Paradigma]
        Estilo de organizar o raciocínio e o código. Linguagens modernas são \textbf{multiparadigma}: Kotlin e Dart misturam imperativo, objetos e funcional.
      \end{definicao}
      \begin{dica}
        Use o paradigma que deixa cada parte mais clara: objetos para modelar o domínio, funções para transformar dados, declarativo para descrever a interface.
      \end{dica}
    \end{column}
  \end{columns}
  \begin{historia}
    Fortran (1957) era imperativo; Lisp (1958), funcional; Simula (1967) e Smalltalk (anos 1970) criaram a orientação a objetos. Kotlin (2011, JetBrains) e Dart (2011, Google) herdaram ideias de todos.
  \end{historia}
\end{frame}

\begin{frame}{Sistemas de tipos}
  \begin{columns}[T]
    \begin{column}{0.5\textwidth}
      \begin{definicao}[Tipo]
        Conjunto de valores possíveis e das operações permitidas sobre eles. Somar dois \cod{Int} faz sentido; somar um \cod{Int} a uma \cod{Bancada}, não.
      \end{definicao}
      \scriptsize
      \begin{tabularx}{\linewidth}{@{}l L@{}}
        \toprule
        \cabecalho{Classificação} & \cabecalho{Significado} \\
        \midrule
        Estática & tipos verificados na compilação (Kotlin, Dart, Java) \\
        Dinâmica & tipos verificados na execução (Python, JavaScript) \\
        Forte & conversões perigosas são proibidas ou explícitas \\
        Inferência & o compilador deduz o tipo (\cod{val x = 3}) \\
        \bottomrule
      \end{tabularx}
    \end{column}
    \begin{column}{0.48\textwidth}
      \begin{conceito}[Por que tipagem estática ajuda]
        Muitos erros aparecem \textbf{antes} de o app rodar; a IDE completa código e refatora com segurança; o tipo documenta a intenção.
      \end{conceito}
      \begin{pordentro}[Solidez (\emph{soundness})]
        Um sistema de tipos é sólido quando garante que, se o programa compila, uma variável nunca terá em execução um valor de outro tipo. O Dart 3 é sólido, inclusive quanto a nulos.
      \end{pordentro}
    \end{column}
  \end{columns}
\end{frame}

\begin{frame}{O ``erro de um bilhão de dólares'': o nulo}
  \begin{columns}[T]
    \begin{column}{0.5\textwidth}
      \begin{historia}
        Tony Hoare introduziu a referência nula na linguagem ALGOL W em 1965. Décadas depois, chamou-a de seu ``erro de um bilhão de dólares'', pelo custo das falhas que ela causou.
      \end{historia}
      \begin{definicao}[\emph{Null safety}]
        O sistema de tipos separa tipos que \textbf{podem} ser nulos (\cod{String?}) dos que \textbf{não podem} (\cod{String}). Usar um valor possivelmente nulo sem tratar vira erro de compilação, não falha em produção.
      \end{definicao}
    \end{column}
    \begin{column}{0.48\textwidth}
      \centering
      \begin{tikzpicture}[node distance=3mm,every node/.style={font=\scriptsize,align=center}]
        \node[caixa,minimum width=2.2cm] (a) {\cod{String?}\\ \tiny pode ser nulo};
        \node[caixapetroleo,right=8mm of a,minimum width=2.2cm] (b) {\cod{String}\\ \tiny nunca nulo};
        \draw[seta] (a) -- node[rotulo,above]{\tiny verificação} (b);
        \node[below=3mm of a,font=\tiny,text width=2.2cm,align=center] {\cod{if (x != null)}\\ \cod{x ?: padrão}};
      \end{tikzpicture}
      \begin{pordentro}[Análise de fluxo]
        Depois de um \cod{if (x != null)}, o compilador sabe que \cod{x} não é nulo naquele trecho e permite usá-lo diretamente (\emph{smart cast} no Kotlin, promoção de tipo no Dart).
      \end{pordentro}
      \begin{atencao}
        O operador \cod{!!} (Kotlin) ou \cod{!} (Dart) desliga a proteção: só use quando tiver certeza.
      \end{atencao}
    \end{column}
  \end{columns}
\end{frame}

\begin{frame}{Os pilares da orientação a objetos}
  \begin{columns}[T]
    \begin{column}{0.54\textwidth}
      \scriptsize
      \begin{tabularx}{\linewidth}{@{}l L@{}}
        \toprule
        \cabecalho{Pilar} & \cabecalho{Significado e exemplo} \\
        \midrule
        Encapsulamento & esconder o estado interno e expor operações: a \cod{Reserva} só muda pelo método \cod{cancelar()} \\
        Abstração & mostrar o essencial e ocultar detalhes: \cod{FonteLeituras} sem dizer se é BLE ou MQTT \\
        Herança & uma classe especializa outra: \cod{SensorTemperatura} herda de \cod{Sensor} \\
        Polimorfismo & o mesmo comando, comportamentos diferentes: \cod{fonte.leituras()} em qualquer implementação \\
        \bottomrule
      \end{tabularx}
    \end{column}
    \begin{column}{0.44\textwidth}
      \begin{definicao}[Objeto e classe]
        A \textbf{classe} é o molde (atributos e métodos); o \textbf{objeto} é uma instância concreta, com estado próprio.
      \end{definicao}
      \begin{dica}[Composição em vez de herança]
        Prefira montar objetos com outros objetos (``tem um'') a criar hierarquias profundas (``é um''): mudanças ficam localizadas e o código, mais flexível.
      \end{dica}
    \end{column}
  \end{columns}
  \begin{analogia}
    Uma tomada é uma abstração: qualquer aparelho com o plugue certo funciona (polimorfismo), e ninguém precisa conhecer a fiação dentro da parede (encapsulamento).
  \end{analogia}
\end{frame}

\begin{frame}{Programação funcional}
  \begin{columns}[T]
    \begin{column}{0.5\textwidth}
      \begin{definicao}[Função pura]
        Para a mesma entrada, devolve sempre a mesma saída e não causa efeitos colaterais (não altera variáveis externas, não grava em disco, não acessa a rede).
      \end{definicao}
      \begin{definicao}[Funções de primeira classe e de ordem superior]
        Funções são valores: podem ser guardadas em variáveis, passadas como argumento e devolvidas. Uma função que recebe ou devolve outra é de \textbf{ordem superior} (\cod{map}, \cod{filter}).
      \end{definicao}
    \end{column}
    \begin{column}{0.48\textwidth}
      \centering
      \begin{tikzpicture}[node distance=4mm,every node/.style={font=\scriptsize,align=center}]
        \node[caixa,minimum width=1.9cm] (l) {leituras};
        \node[caixapetroleo,right=of l,minimum width=1.6cm] (f) {\cod{filter}\\ \tiny > 30\,°C};
        \node[caixapetroleo,right=of f,minimum width=1.6cm] (m) {\cod{map}\\ \tiny para alerta};
        \node[caixaambar,below=4mm of m,minimum width=1.6cm] (r) {alertas};
        \draw[seta] (l) -- (f); \draw[seta] (f) -- (m); \draw[seta] (m) -- (r);
      \end{tikzpicture}
      \begin{conceito}[Por que importa no mobile]
        Funções puras são fáceis de testar e seguras entre \emph{threads}; transformar listas com \cod{map}/\cod{filter} é mais legível que laços; interfaces declarativas tratam a tela como uma função do estado (Módulo 05).
      \end{conceito}
    \end{column}
  \end{columns}
\end{frame}

\begin{frame}{Memória: pilha, \emph{heap} e coleta de lixo}
  \begin{columns}[c]
    \begin{column}{0.46\textwidth}
      \centering
      \begin{tikzpicture}[every node/.style={font=\tiny}]
        \node[font=\scriptsize\bfseries,text=azulpetroleo] at (0.9,2.7) {Pilha};
        \foreach \y/\t in {2.2/main(),1.7/carregar(),1.2/calcular()} \node[caixa,minimum width=1.8cm] at (0.9,\y) {\t};
        \node[font=\scriptsize\bfseries,text=azulpetroleo] at (3.7,2.7) {\emph{Heap}};
        \draw[cinzamedio,rounded corners] (2.6,0.3) rectangle (4.8,2.45);
        \node[caixapetroleo,minimum width=1.3cm] (o1) at (3.3,1.95) {Bancada};
        \node[caixapetroleo,minimum width=1.3cm] (o2) at (4.1,1.2) {Lista};
        \node[caixa,minimum width=1.3cm,fill=cinzaclaro] (o3) at (3.3,0.65) {sem ref.};
        \draw[seta] (1.8,2.2) -- (o1); \draw[seta] (1.8,1.7) -- (o2); \draw[seta] (o2) -- (o1);
        \node[font=\tiny,text=vermelhoerro] at (3.7,0.2) {coletado};
      \end{tikzpicture}
    \end{column}
    \begin{column}{0.52\textwidth}
      \begin{definicao}[Pilha e \emph{heap}]
        A \textbf{pilha} guarda as chamadas de função e suas variáveis locais, liberadas ao retornar. O \textbf{\emph{heap}} guarda os objetos, que vivem enquanto houver referências a eles.
      \end{definicao}
      \begin{definicao}[Coleta de lixo]
        O \emph{runtime} encontra objetos que não são mais alcançáveis a partir das referências vivas e libera sua memória automaticamente.
      \end{definicao}
      \begin{pordentro}[Gerações]
        A maioria dos objetos morre jovem. Por isso o ART e a VM do Dart separam objetos novos e antigos e coletam os novos com frequência e rapidez.
      \end{pordentro}
    \end{column}
  \end{columns}
\end{frame}

\begin{frame}{Imutabilidade e igualdade}
  \begin{columns}[T]
    \begin{column}{0.5\textwidth}
      \begin{definicao}[Imutável]
        Objeto que não muda depois de criado. Para ``alterar'', cria-se uma cópia com o campo novo (\cod{copy} no Kotlin, \cod{copyWith} no Dart).
      \end{definicao}
      \begin{conceito}[Vantagens]
        \begin{itemize}
          \item Ninguém altera o objeto ``pelas costas''.
          \item Seguro para compartilhar entre \emph{threads} e \emph{isolates}.
          \item Comparar o antes e o depois de um estado fica trivial (Módulo 07).
        \end{itemize}
      \end{conceito}
    \end{column}
    \begin{column}{0.48\textwidth}
      \begin{definicao}[Igualdade de valor e de identidade]
        \textbf{Identidade}: as duas variáveis apontam para o \textbf{mesmo} objeto. \textbf{Valor}: objetos diferentes com o mesmo conteúdo. \cod{data class} (Kotlin) e \emph{records} (Dart) comparam por valor.
      \end{definicao}
      \begin{erro}[Lista mutável compartilhada]
        Dois componentes recebem a mesma lista; um a altera e o outro mostra dados que não esperava. Exponha listas somente leitura.
      \end{erro}
    \end{column}
  \end{columns}
\end{frame}
